% Generated draft: requires mathematical and source review.
\documentclass[11pt]{article}
\usepackage[margin=25mm]{geometry}
\usepackage{fontspec}
\setmainfont{Latin Modern Roman}
\usepackage{amsmath,amssymb,mathtools,unicode-math}
\setmathfont{Latin Modern Math}
\usepackage{enumitem,needspace,xurl,hyperref}
\setlength{\parindent}{0pt}
\setlength{\parskip}{5pt}
\setlength{\emergencystretch}{3em}
\newcommand{\N}{\mathbb{N}}
\newcommand{\Z}{\mathbb{Z}}
\newcommand{\Q}{\mathbb{Q}}
\newcommand{\R}{\mathbb{R}}
\newcommand{\C}{\mathbb{C}}
\newcommand{\F}{\mathbb{F}}
\newcommand{\A}{\mathbb{A}}
\newcommand{\PP}{\mathbb{P}}
\newcommand{\E}{\mathbb{E}}
\newcommand{\eps}{\varepsilon}
\newcommand{\dd}{\,\mathrm{d}}
\begin{document}
\section*{UnsolvedMath ID: 30002351}
\textbf{Source problem code:} catalogue entry
\subsection{Definitions and mathematical statement}
For context, the usual degree-two norm-residue formulation of the Merkurjev--Suslin theorem is the isomorphism
\[
K_2^M(F)/mK_2^M(F)\xrightarrow{\ \sim\ }H^2(G_F,\mu_m^{\otimes 2}),
\qquad \{a,b\}\longmapsto (a)\smile(b),
\]
for a field \(F\) and an integer \(m\geq 2\) invertible in \(F\). Here \(G_F=\operatorname{Gal}(F^{\mathrm{sep}}/F)\), \(\mu_m\) is the Galois module of \(m\)-th roots of unity, cohomology is continuous Galois cohomology, and \((a)\in H^1(G_F,\mu_m)\) is the Kummer class of \(a\in F^\times\). The group \(K_2^M(F)\) is generated by bilinear symbols \(\{a,b\}\) subject to the Steinberg relations \(\{a,1-a\}=0\) for \(a\neq 0,1\). This recalls the conventional theorem mentioned in the question; it does not specify the intended higher analogue.

For a standard degree-one convention for Massey products, let \(A^\bullet\) be a differential graded algebra with differential \(d\), and let \(\alpha_1,\ldots,\alpha_n\in H^1(A^\bullet)\), where \(n\geq 2\). A defining system consists of elements \(a_{ij}\in A^1\), indexed by \(1\leq i<j\leq n+1\) with \((i,j)\neq(1,n+1)\), such that \(a_{i,i+1}\) represents \(\alpha_i\) and
\[
d a_{ij}=\sum_{i<k<j}a_{ik}a_{kj}.
\]
With this sign convention, the Massey product is the set
\[
\langle\alpha_1,\ldots,\alpha_n\rangle
=
\left\{
\left[\sum_{k=2}^{n}a_{1k}a_{k,n+1}\right]
:\ (a_{ij})\text{ is a defining system}
\right\}\subseteq H^2(A^\bullet).
\]
It is defined when a defining system exists, and it vanishes when it contains zero. For \(n=2\), it is the singleton containing \(\alpha_1\smile\alpha_2\). These definitions explain one conventional setting; the supplied record does not specify the degrees, coefficients, or differential graded algebra intended in its question.

\textbf{Source question.} Can the Merkurjev--Suslin theorem be generalized from the twofold cup product to \(n\)-fold Massey products?

\textbf{Unresolved formulation.} The supplied question does not define what constitutes such a generalization. In particular, it supplies no higher symbol group, comparison map, treatment of the multiple possible values of a Massey product, or required conclusion. It also does not specify the fields, coefficient modules, root-of-unity hypotheses, input degrees, or quantification over defining systems. Thus the record supports the research question above, but does not determine a precise universally quantified conjecture. The general vanishing question mentioned in the supplied background must not be identified with a Merkurjev--Suslin-type comparison theorem without additional source evidence.
\subsection{Short English statement}
The question asks for an analogue of the Merkurjev--Suslin theorem involving higher Massey products. Such products can have several values and require defining systems. The supplied record does not specify the proposed analogue or its hypotheses, so a unique precise conjecture cannot be extracted. Its open-status and literature claims are supplied dated curation requiring independent review.
\subsection{Sources}
\begin{itemize}
\item Local UnsolvedMath registry, supplied record with permanent ID \(30002351\), title \emph{Merkurjev--Suslin Theory for Higher Massey Products}, exact source problem code \texttt{catalogue entry}, in the collection \emph{Oberwolfach Reports Open Problems}. The record gives proposed year 2013 and names no proposer.
\item Supplied report citation: \emph{The Arithmetic of Fields (2014)}, Oberwolfach Reports, DOI 10.4171/owr/2013/30, pp.~1771--1818. \url{https://doi.org/10.4171/owr/2013/30}. The record describes these reports as documenting mathematical workshops and their open-problem discussions. No problem-specific page or original passage is supplied.
\item Efrat--Matzri, \emph{Triple Massey products and absolute Galois groups}. \url{https://arxiv.org/abs/1412.7265}.
\item Guillot--Minac, \emph{Extensions of unipotent groups and Massey products}. \url{https://arxiv.org/abs/1711.07711}.
\item Reference supplied under the label \emph{Advances in Mathematics publication}. \url{https://doi.org/10.1016/j.aim.2019.106748}. The record supplies no further bibliographic identification for this entry.
\item Literature assessment dated exactly \texttt{2026-08-21}: the supplied curation labels the problem open and states that triple Massey products vanish for absolute Galois groups, that substantial fourfold and higher special cases are known, and that the requested all-\(n\) generalization and general Massey-vanishing and Merkurjev--Suslin-type statement remain open. These are supplied dated curation claims, not independently verified mathematical facts or a new literature assessment. The record itself requests independent verification of time-sensitive status and cited claims.
\item The registry labels the statement \texttt{corrected\_verified} and reports manual semantic, typographic, and LaTeX review of a meaning-preserving cleanup. This supplied editorial label does not resolve the mathematical formulation gaps or establish independent source verification.
\end{itemize}
\end{document}
